\chapter{Bestyrelsesgrupperinger}

\section{\CERM{}lauget}

Den mest truttende og båttende underforening på \TKET består af siddende og tidligere ceremonimestre og viceceremonimestre. I fordums tid har lauget været i åben væbnet indre konflikt mellem de to under-underforeninger ``elite-\CERM{}lauget'' og ``elite-elite\CERM{}lauget'', som bestod af henholdsvist ceremonimestre og viceceremonimestre. Striden endte uafgjort, og underforeningen har siden fungeret upåklageligt og meget effektivt.\par

\CERM{}laugets fornemste opgave består i den løbende vedligeholdelse og tilbedelse af vagtbrien Brian. Hver torsdag ved fuldmåne iklæder underforeningen sig rituelle kutter, tænder rituelle fakler, og fodrer for alle medlemmers åsyn vagtbrien med brie og top under fuldmånens skær.\par

\CERM{}lauget er af denne og andre grunde den vigtigste og mest betydningsfulde samling af individer på \TKET --- nej, i Aarhus!\par

Af andre vigtige (dog ikke så vigtige som Brianfodringen) opgaver og aktiviteter, kan også nævnes den årlige syvogfyrrende dag. Her mødes \CERM{}lauget på årets syvogfyrresinstyvende dag på året (typisk 16. februar, alt efter skudåret) og nyder iskolde Carlsberg 47, mens de mange sange fra indeværende samt forgangne syvogfyrrende dage synges i smuk harmoni. Således bliver arrangementet år for år én sang længere.\par

Den første onsdag i maj klokken 12.00 tages kampen op med de fandenivoldske sirener, der til evig tid forpester det danske godtfolks øregange. Her gør \CERM{}lauget hvad de kan, for at overdøve spektaklet med rene, klare toner fra deres mange horn og øvrige blæseinstrumenter. Hermed ser man tydeligt underforeningens altruisme og velgørende sjæl, når fjenden hvert år besejres!\par

\CERM{}laugets medlemmer genkendes lettest på deres flagrende kutter (hvoraf Ute Hahn har syet fire) og lållende røster (oftest ``The Imperial March'', men nogle gange ``Bubbibjørnene''). Derudover kan yngste (i anciennitet) \CERM og \VC genkendes på deres funklende relikvier. \CERM bærer til dagligt sin flotte kappe, der er sort udenpå og violet indeni, når den ikke er violet udenpå og sort indeni. Derudover bærer hen et flot, poleret jagthorn og en ældgammel \CERM{}s hat med meget lang skygge. Til særligt formelle formål bæres \CERM{}kæden, hvorpå \CERM{}s navn er graveret. \VC bærer ingen kappe, men ligeledes horn. \VC{}-hornet er småt, da \VC{}'er forstår at få meget ud af lidt. Derudover har \VC bræt, som hen bærer om halsen med glæde og stolthed, sammen med sit \VC{}-mærke, som typisk er et omvendt CERES midt på bugen.\par


\section{Formjuntaen}
\tkpicture[\textwidth]{../img/formjuntaen}
\section{Næstformationen}
\tkpicture[0.8\linewidth]{../img/næstformationen}

\section{Kasserernes Efterretningstjeneste (KET)}

KET er organisationen af nuværende og tidligere \KASS{}erere og INKAssatorer for \TKET. KET administrerer \TKETs likvide tequila, flydende midler, faste ejendomme og øvrige nationale og oversøiske besiddelser.

Det er KETs naturlige opgave at opspore gamle og nye skyldnere og udvirke at de ophører med at være dette. Til dette formål har \KASS en kugleramme til at udregne alle former for gæld, moms og renter i en vilkårlig valuta. Desuden har KET placeret agenter på strategiske poster i både ind- og udland, og allieret sig med både Experian og HA Inkasso af hensyn til de mere genstridige skyldnere. Hvis disse midler ikke er tilstrækkelige til at få folk til at erkende gælden og erlægge det skyldige beløb, råder KET over et antal rokokokøller. Synet af disse i en kyndig INKA’s hånd bringer som oftest folk på bedre tanker.

Når KET mødes for at diskutere \TKETs midler, den danske og internationale økonomi, eller hvordan vi får indført planøkonomi med \KASS og INKA i førersædet, kræver det hovedet. Udfordringer som \TKET og resten af verden har stødt på de seneste 10 år, har ikke kun påvirket økonomien men også hygiejnen. Vi har alle lært, at det er vigtigt at spritte af efter brug. Derfor vælger vi selvfølgelig i KET at spritte af med den ondeste sprit,Tequila! Dette sørger for at hjerne bliver godt renset, endda så godt at man glemmer hvordan man lavede de beregninger til KET mødet, så man bliver nødt til at regne på dem igen for at dobbelttjekke deres validitet.

En af de fornemmeste opgaver i KET er at sikre den bedste forvaltning af \TKETs midler. For at sikre dette kan det være nødvendigt at modarbejde de øvrige bestyrelsesmedlemmer for at stoppe deres ødslen med de begrænsede midler, som er til rådighed i den daglige drift. 
Ud over ovenstående sørger KET for at drikke øl og støtte ethvert andet formål,
som på sin vis bidrager til at forbedre \TKETs økonomi.


KET ønsker \TKET tillykke med de 70 år. Selvom \TKETs økonomi fik en lille virus i løbet af de sidste 10 år, tror vi på at \TKETs øltørst vil holde både den danske og \TKETs økonomi stærk og i øvrigt mener vi at \TKET vil bestå!

\section{Spiril}

Spirillen er et lille, gult dyr med sorte pletter. Det lever i Palombias jungle og har efter Disneys overtagelse fået både unger og en mage. Det har en utrolig lang hale, som den bruger til hvad den har lyst, trods den er et amfibium.

Dette har på ingen måde noget med André Franquin, der holdt op med ikke at lade være med at tegne serien, da han undlod at afstå fra at dø, eller TÅGEKAMMERET at gøre. Hvis det havde det, ville SPIRIL i hvert fald ikke være en forkortelse for SEKR’r og PR’r I Ret Intelligent Loge. Og skulle det have været tilfældet, holdt denne loge absolut ikke møde nogensinde.

Hvis ikke den under ingen omstændigheder ej nogensinde havde undladt at holde et møde, ville der i hvert fald ikke efterfølgende komme udemokratiske opkrævninger fra Post Danmark for første gang. 

Som nogen engang skrev: Vi ved ikke, hvorfor vi ikke skriver dette

\begin{flushright}
{/Inge N.}
\end{flushright}
\section{FU’ernes CentralKomité}

\textbf{FUCK}, eller \textbf{FU}’ernes CentralKomité, er foreningen af gamle og nuværende FU’er. Foreningen blev oprettet under den stiftende generalforsamling den 14. september 1988 og har spredt glæde og sammenhold blandt FU’er lige siden.

I \textbf{FUCK}s vedtægter finder man, at formålet med foreningen er: “indimellem at mødes under hyggelige former, over adskillige kasser øl – til sanseløs druk.” I selvsamme vedtægter findes ligeledes, at foreningen har en formand, som er FULD (FU-LeDer). Af samme grund tager OFULD (Officiel FU-LeDer) sig af den daglige arbejdsbyrde. Kun OFULD kender til FULDs identitet, så denne ikke kan blive modarbejdet i at varetage FU’ernes interesser. Først til TÅGEKAMMERETs generalforsamling giver FULD sig til kende og holder en tale.

Hvert år afholdes i \textbf{FUCK}-regi en generalforsamling og en julefrokost. Herudover kan OFULD (eller andre medlemmer) indkalde til andre arrangementer, såfremt denne finder det nødvendigt (eller i det mindste sjovt).

Alle \textbf{FUCK}s aktiviteter nedfældes i referatet af FUSK (FUSeKretær), således at intet bliver glemt (selvom ting som ellers ikke var husket kan være svære at få til at give mening).

De FU’er, som bliver valgt til TÅGEKAMMERETs bestyrelse, bliver midlertidigt udkastet af foreningen og kan først genoptages, når de forlader førnævnte bestyrelse. Selv når dette er opfyldt, må de underskrive en løjalitetserklæring, og udføre visse opgaver (som er nærmere specificeret i den enkeltes løjalitetserklæring) for at råde bod, i vor vel og såfremt at de så ganske ønsker genoptagelse i den fantastiske og uforgængelige forening \textbf{FUCK}{\tiny{\rule[0pt]{0.9em}{0.9em}$\longleftarrow$\,gotisk punktum}}

\section{Brev til Ålborg flyvestation}

Det er ikke altid lige let at finde folk som kan og vil stille en bil til rådighed ved
HSTR. I sådanne situationer er det oplagt at henvende sig til militæret. Det var i
hvert fald hvad ceremonimesteren gjorde i 1982.

\tkpicture[0.9\textwidth]{../img/aalborg_lufthavn}

\vspace{-130pt}

\noindent
Kommandantens meget imødekommende svar kan ses i \TKETs arkiv.
